\documentclass{beamer}
\usepackage{ctex, hyperref}
\usepackage[T1]{fontenc}

% other packages
\usepackage{latexsym,xcolor,multicol,booktabs,calligra}
\usepackage{amsmath,amssymb,amsfonts}
\usepackage{graphicx,pstricks,listings,stackengine}
\usepackage{float}
\usepackage{subfig}
\usepackage{algorithmic}
\usepackage{algorithm}
\usepackage{multirow}

\author{学生：毛周祥~~~指导老师：单冯}
\title{无标注预训练方法在视觉语言模型的应用}
\subtitle{科学研究方法--课程作业}
\institute{东南大学计算机拔尖培养基地}
\date{2025 年 9 月}
\usepackage{SEU}

% defs
\def\cmd#1{\texttt{\color{red}\footnotesize $\backslash$#1}}
\def\env#1{\texttt{\color{blue}\footnotesize #1}}
\definecolor{seu-green}{rgb}{0.345,0.459,0.345}      % #587558
\definecolor{seu-yellow}{rgb}{0.992,0.816,0}         % #fdd000
\definecolor{seu-black}{rgb}{0.137,0.094,0.082}      % #231815
\definecolor{deepblue}{rgb}{0,0,0.5}
\definecolor{deepred}{rgb}{0.6,0,0}
\setbeamerfont{normal text}{size=\Large}
\setbeamerfont{sub normal text}{size=\large}

% 导言区
\makeatletter
\newenvironment{Litemize}{%
  \par\begingroup
  \def\beamer@setuplines{\baselineskip=36pt\relax}% 行距
  \usebeamerfont{itemize item}% 让 beamer 保留标签设置
  \Large\linespread{1}\selectfont      % 字号 + 行距
  \begin{itemize}[<+->]
}{%
  \end{itemize}%
  \endgroup
}

\newcommand{\subtit}[1]{\\[2pt]  % 换行+微小垂直间隙
   {\normalsize\color{gray!70} #1 \par}}   % 灰色、小一号


\newenvironment{litemize}{%
  \par\begingroup
  \def\beamer@setuplines{\baselineskip=12pt\relax}%
  \usebeamerfont{itemize item}%
  \large\linespread{1.25}\selectfont
  \begin{itemize}[<+->]
}{%
  \end{itemize}%
  \endgroup
}

\newenvironment{litemize-p}{%
  \par\begingroup
  \def\beamer@setuplines{\baselineskip=12pt\relax}%
  \usebeamerfont{itemize item}%
  \large\linespread{1.25}\selectfont
  \begin{itemize}[]
}{%
  \end{itemize}%
  \endgroup
}
\makeatother

\lstset{
    basicstyle=\ttfamily\small,
    keywordstyle=\bfseries\color{deepblue},
    emphstyle=\ttfamily\color{deepred},    % Custom highlighting style
    stringstyle=\color{seu-green},
    numbers=left,
    numberstyle=\small\color{seu-black},
    rulesepcolor=\color{red!20!green!20!blue!20},
    frame=shadowbox,
}


\begin{document}


\kaishu
\begin{frame}
    \titlepage
    \begin{figure}[htpb]
        \begin{center}
            \includegraphics[width=0.33\linewidth]{pic/SEU-09J.png}
        \end{center}
    \end{figure}
\end{frame}

\begin{frame}
    \tableofcontents[sectionstyle=show,subsectionstyle=show/shaded/hide,subsubsectionstyle=show/shaded/hide]
\end{frame}

\section{研究背景}

\begin{frame}


    \frametitle{为什么需要视觉语言模型 (VLM)？}
    \begin{columns}
        \column{0.5\textwidth}
        \begin{Litemize}
            \item 一句话定义 VLM
            \item 它能做什么？ 
                \begin{litemize-p}
                    \item 图文配对理解
                    \item 视觉问答（VQA）
                    \item 图像字幕生成
                \end{litemize-p}
        \end{Litemize}
        \column{0.5\textwidth}
        \begin{figure}[htpb]
            \begin{center}
                \includegraphics<3->[width=1\linewidth]{pic/VLM-usage.png}
            \end{center}
        \end{figure}
    \end{columns}
    
\end{frame}

\begin{frame}
    \frametitle{无标注预训练方法的 Motivation}
    \begin{columns}
        \column{0.5\textwidth}
        \begin{Litemize}
            \item 现有方法的局限
            \begin{litemize-p}
                \item 依赖昂贵的有标注数据
                \item 难以扩展到新任务
                \item 图像区域与文本词汇难以精细对齐
            \end{litemize-p}
            \item NLP 领域的成功经验
            \begin{litemize-p}
                \item 从 Word2Vec 到 GPT
                \item 自监督预训练
            \end{litemize-p}      
        \end{Litemize}
        \column{0.5\textwidth}
        \begin{figure}[htpb]
            \begin{center}
                \includegraphics<1->[width=1\linewidth]{pic/why-pretraining-1.png}
            \end{center}
        \end{figure}
        \begin{figure}[htpb]
            \begin{center}
                \includegraphics<5->[width=1\linewidth]{pic/why-pretraining-2.png}
            \end{center}
        \end{figure}
    \end{columns}
\end{frame}

\begin{frame}
    \frametitle{大规模无标注预训练 \\
    ~~~~~~----- 观千剑而后识器，操千曲而后晓声
    }
    \begin{columns}
        \column{0.5\textwidth}
        \begin{Litemize}
            \item 突破标注瓶颈
            \item \color{seu-yellow}零样本泛化\color{black}涌现
            \item 统一表示空间
            \item 通用下游任务范式
        \end{Litemize}
        \column{0.5\textwidth}
        \begin{figure}[htpb]
            \begin{center}
                \includegraphics<4->[width=1\linewidth]{pic/VLM-pretraining-workflow.png}
            \end{center}
        \end{figure}
    \end{columns}
\end{frame}

\section{CLIP}
\begin{frame}
    \frametitle{CLIP (2021)}
    \begin{columns}
        \column{0.6\textwidth}
        \begin{litemize}
            \item 目标：学习图文的联合表示
            \item 方法：对比学习
            \item 数据集：4 亿对图文数据
            \item 结果：多种下游任务表现优异
        \end{litemize}
        \column{0.5\textwidth}
        \begin{figure}[htpb]
            \begin{center}
                \includegraphics<1->[width=1\linewidth]{pic/CLIP-demo.png}
            \end{center}
        \end{figure}
    \end{columns}
\end{frame}

\begin{frame}
    \frametitle{CLIP 的架构}
    \begin{columns}
        \column{0.5\textwidth}
        \begin{litemize}
            \item 图像编码器：ResNet 或 ViT
            \item 文本编码器：Text Transformer
            \item 总结构：双塔架构
        \end{litemize}
        \column{0.5\textwidth}
        \begin{figure}[htpb]
            \begin{center}
                \includegraphics<1>[width=1.2\linewidth]{pic/CLIP-architecture-3.png}
            \end{center}
        \end{figure}
        \begin{figure}[htpb]
            \begin{center}
                \includegraphics<2>[width=1\linewidth]{pic/CLIP-architecture-2.png}
            \end{center}
        \end{figure}
        \begin{figure}[htpb]
            \begin{center}
                \includegraphics<3>[width=1\linewidth]{pic/CLIP-architecture-1.png}
            \end{center}
        \end{figure}
    \end{columns}
\end{frame}
\begin{frame}
    \frametitle{CLIP 的训练——对比学习}
    \begin{columns}
        \column{0.5\textwidth}
        \begin{litemize}
            \item 训练方法：图文对比 (ITC)
            \item 优化函数：对比损失
            \item 为什么有效？
            \begin{litemize-p}
                \item 损失函数就是考试成绩
                \item 倒逼模型学习一个合理的图文表示
                \item 不相关的推远，相关的拉近
                \item 大规模数据+大模型
            \end{litemize-p}
        \end{litemize}
        \column{0.5\textwidth}
        \begin{figure}[htpb]
            \begin{center}
                \includegraphics<1->[width=1\linewidth]{pic/CLIP-ITC.png}
            \end{center}
        \end{figure}
    \end{columns}
\end{frame}
\begin{frame}
    \frametitle{CLIP 的训练——队列}
    \begin{columns}
        \column{0.5\textwidth}
        \begin{litemize}
            \item 四亿张图需要两两对比？
            \item 批次大小受限，负例不足。
            \item 解决方案：队列 (Queue)
        \end{litemize}
        \begin{figure}
            \begin{center}
                \includegraphics<4->[width=0.7\linewidth]{pic/GM-contrast.png}
            \end{center}
        \end{figure}
        \column{0.5\textwidth}
        \begin{figure}[htpb]
            \begin{center}
                \includegraphics<1->[width=1\linewidth]{pic/CLIP-ITC.png}
            \end{center}
        \end{figure}
    \end{columns}
\end{frame}
\begin{frame}
    \frametitle{CLIP 的下游使用}
    \begin{columns}
        \column{0.5\textwidth}
        \begin{litemize}
            \item ZeroShot 图像分类。
            \item 作为其他任务的编码器。\\
            eg. Stable Diffusion
            \item 为什么能迁移？
            \begin{litemize-p}
                \item 语言本身就是监督信号
                \item 只有学习到如何表示图片特征，才能完成预训练任务 \\
                只有学会了课上讲的内容，才能独立完成实验作业
            \end{litemize-p}
        \end{litemize}

        \column{0.5\textwidth}
        \begin{figure}[htpb]
            \begin{center}
                \includegraphics<1->[width=1\linewidth]{pic/CLIP-zeroshot.png}
            \end{center}
        \end{figure}
    \end{columns}
\end{frame}

\section{ALBEF}

\begin{frame}
    \frametitle{CLIP 的缺陷}
    \begin{columns}
        \column{0.5\textwidth}
        \begin{litemize}
            \item 细粒度对齐能力不足
            \begin{litemize-p}
                \item 训练时只针对图文整体
                \item 无法精确匹配图像区域与文本词汇
                \item 难以处理复杂场景
            \end{litemize-p}
            \item 数据集质量参差不齐
            \begin{litemize-p}
                \item 网络爬取，噪声多
                \item 不准确标签极大影响模型学习
            \end{litemize-p}
        \end{litemize}
        \column{0.5\textwidth}
        \begin{figure}[htpb]
            \begin{center}
                \includegraphics<5->[width=1\linewidth]{pic/web-noise.png}
            \end{center}
        \end{figure}
    \end{columns}
\end{frame}
\begin{frame}
    \frametitle{ALBEF (2021)}
    \begin{itemize}
        \item 第一部分结构与 CLIP 相同
        \item 增加了跨模态融合模块（Cross Attention）
        \item 增加了 ITM 和 MLM 的预训练任务，增强 Text-Image 交互。
    \end{itemize}
    \begin{figure}
        \begin{center}
            \includegraphics<1->[width=0.8\linewidth]{pic/ALBEF.png}
        \end{center}
    \end{figure}
\end{frame}
\begin{frame}
    \frametitle{ITM + MLM}
    \begin{litemize}
        \item Image Text Matching (ITM)
        \begin{litemize-p}
            \item ITC 只是大致对齐信息，ITM 可以让模型学会图文交互。
            \item hard-negative 采样：通过细微差异制造难负例，倒逼模型学习交互
            \item 在多模态输出层加一个 FC Layer 输出概率分布，然后反向传播
        \end{litemize-p}
        \item Masked Language Modeling (MLM)
        \begin{itemize}
            \item 随机遮蔽文本中的部分词汇
            \item 根据剩余信息，预测遮蔽词
            \item 要准确预测，必须能够关注到图像中相关区域，并理解图文交互逻辑
        \end{itemize}
    \end{litemize}
\end{frame}
\begin{frame}
    \frametitle{动量蒸馏 (MoD)}
        \begin{litemize}
            \item 动量模型提供软标签
            \begin{figure}[htpb]
                \begin{center}
                    \includegraphics<1>[width=1\linewidth]{pic/pseudo-targets.png}
                \end{center}
            \end{figure}
            \item 加权 KL 散度做损失函数
            \begin{math}
            \mathcal{L} = \mathbb{E}_{(I,\hat{T})\sim D_{\mathcal{H}}}
            \bigl[-\log p_{\text{msk}}(I,\hat{T})\bigr]
            \end{math} \\
            \begin{math}
            \mathcal{L}' = (1-\alpha) \mathcal{L}
            + \alpha\,\mathbb{E}_{(I,\hat{T})\sim D_{\mathcal{K}\!\mathcal{L}}}
            \bigl[\mathrm{KL}\!\bigl(q_{\text{msk}}(I,\hat{T})\big\|p_{\text{msk}}(I,\hat{T})\bigr)\bigr]
            \end{math}
            \item 天然降噪不过罚噪声样本
        \end{litemize}
        
\end{frame}


\section{BLIP}
\begin{frame}
    \frametitle{BLIP (2022)}
    \begin{itemize}
        \item ALBEF 并没有从根本上解决数据噪声问题
        \item BLIP 提出\color{seu-yellow}数据自举 (Bootstrapping) \color{black}方法
        \item 为了做出改进，BLIP 将 ALBEF 的 MLM 任务更换为了 LM (语言建模) 任务，以便完成自举
    \end{itemize}
    \begin{figure}
        \begin{center}
            \includegraphics<1->[width=0.8\linewidth]{pic/BLIP-architecture.png}
        \end{center}
    \end{figure}
\end{frame}
\begin{frame}
    \frametitle{LM 建模}
    \begin{itemize}
        \item Causal Self-Attention，单向预测
        \item Cross Attention，融合图片信息，图文交互
        \item 生成式损失，生成下一个词的分布，交叉熵损失
    \end{itemize}
    \begin{figure}
        \begin{center}
            \includegraphics<1->[width=0.8\linewidth]{pic/BLIP-LM.png}
        \end{center}
    \end{figure}
\end{frame}
\begin{frame}
    \frametitle{数据自举 (Boostraping)}
    \begin{itemize}
        \item 少量有标注数据预训练
        \item 对预训练的模型做不同的 fine-tuning 得到 Filter 和 Captizoner
        \item 用 Captizoner 生成伪标签，用 Filter 过滤，保留能通过过滤的所有数据对
    \end{itemize}
    \begin{figure}
        \begin{center}
            \includegraphics<1->[width=1\linewidth]{pic/BLIP-bootstrapping.png}
        \end{center}
    \end{figure}
\end{frame}
\begin{frame}
    \frametitle{性能对比}
    \begin{itemize}
        \item 下游任务：图文检索 (TR) 和图像检索 (IR)
        \item 数据集：Flickr30K
        \item 评价指标：R1 召回率 (一次性正确率)
    \end{itemize}
    \begin{figure}
        \begin{center}
            \includegraphics<1>[width=0.8\linewidth]{pic/compare-acc.png}
        \end{center}
    \end{figure}
    \vspace{-2em}
    \begin{figure}
        \begin{center}
            \includegraphics<2>[width=0.8\linewidth]{pic/compare-err.png}
        \end{center}
    \end{figure}
\end{frame}
\section{拓展介绍}
\begin{frame}
    \frametitle{下游泛化——prompt tuning (2023)}
\end{frame}
\begin{frame}
    \frametitle{下游泛化——Text as Image (2023)}
\end{frame}
\begin{frame}
    \frametitle{应用生态}
\end{frame}

\section{最后}
\begin{frame}
    \frametitle{总结}
    \begin{itemize}
        \item VLM 及其预训练方法的重要性
        \item 为什么预训练任务有效
        \begin{litemize-p}
            \item 预训练任务是数学考试
            \item 训练目标是考高分、倒逼模型学好特征
            \item 学好特征，更好的迁移到下游任务
        \end{litemize-p}
        \item VLM 的演进
        \begin{litemize-p}
            \item CLIP：对比学习，双塔架构
            \item ALBEF：细粒度对齐尝试，动量蒸馏
            \item BLIP：数据自举
        \end{litemize-p}
    \end{itemize}
\end{frame}

\begin{frame}
    \frametitle{Q \& A}
\end{frame}

\begin{frame}
    \begin{center}
        {\Huge\calligra Thanks!}
    \end{center}
\end{frame}

\end{document}